\documentclass[11pt,a4paper]{article}

% --- Pacotes de Idioma e Codificação ---
\usepackage[utf8]{inputenc}
\usepackage[T1]{fontenc}
\usepackage[portuguese]{babel}

% --- Design e Layout ---
\usepackage[margin=2.5cm]{geometry}
\usepackage{charter} % Fonte profissional e legível
\usepackage{lastpage}
\usepackage{fancyhdr}
\usepackage[table]{xcolor}
\usepackage{tabularx}
\usepackage{booktabs}
\usepackage{xcolor}
\usepackage{enumitem}
\usepackage{amssymb}

% --- Configuração de Cores ---
\definecolor{primaryblue}{RGB}{0, 51, 102}

% --- Cabeçalho e Rodapé ---
\pagestyle{fancy}
\fancyhf{}
\fancyhead[L]{\small \textbf{Registo de Entrevista} | EcoRide Solutions}
\fancyhead[R]{\small \thepage\ de \pageref{LastPage}}

\renewcommand{\headrulewidth}{0.4pt}

% --- Ambiente Personalizado para Transcrição ---
% Facilita a distinção entre Entrevistador (E) e Entrevistado (R)
\newlist{transcription}{description}{1}
\setlist[transcription]{
	labelsep=1em,
	leftmargin=3.5em,
	style=nextline,
	font=\bfseries\color{primaryblue}
}

\newcommand{\pergunta}[1]{\item[P:] \textit{#1}}
\newcommand{\resposta}[1]{\item[R:] #1}

% --- Início do Documento ---
\begin{document}
	
	% Título Principal
	{\noindent\LARGE\bfseries\color{primaryblue} Relatório de Entrevista Individual}
	\\[0.5cm]
	
	% --- Tabela de Metadados ---
	\renewcommand{\arraystretch}{1.5}
	\begin{tabularx}{\textwidth}{@{}|l|X|l|X|@{}}
		\hline
		\rowcolor{gray!10} \textbf{ID Entrevista} & \#002 & \textbf{Data} & 26 de Fevereiro de 2026 \\ \hline
		\textbf{Hora Início} & 11:30 & \textbf{Hora Fim} & 12:10 \\ \hline
		\textbf{Local/Meio} & Sala de Reuniões B & \textbf{Estado} & Finalizada / Transcrita \\ \hline
	\end{tabularx}
	
	\vspace{0.5cm}
	
	% --- Participantes ---
	\begin{tabularx}{\textwidth}{@{}|l|X|@{}}
		\hline
		\rowcolor{gray!10} \textbf{Intervenientes} & \textbf{Cargo / Função} \\ \hline
		\textbf{Entrevistador(a)} & [O Teu Nome] - Consultor de Processos \\ \hline
		\textbf{Entrevistado(a)}  & [Nome do Colaborador] - Assistente Administrativa \\ \hline
	\end{tabularx}
	
	\vspace{0.8cm}
	
	% --- Contexto e Objetivos ---
	\section*{1. Enquadramento e Objetivos}
	\begin{itemize}[leftmargin=*, label=\small\color{primaryblue}$\blacksquare$]
		\item \textbf{Otimização do Tempo de Balcão}: O objetivo é garantir que o registo de entrada seja feito em segundos, evitando filas e garantindo que nenhum dado vital se perde.
		\item \textbf{Melhoria da Comunicação}: As perguntas visam perceber como o software pode atuar como ponte entre a oficina e o cliente, usando a funcionária como pivot central.
		\item \textbf{Redução de Erros e Conflitos}: Foca em garantir que a informação de pagamento e o estado do equipamento à entrada sejam incontestáveis, protegendo a funcionária de reclamações.
	\end{itemize}
	
	\hrule % Linha separadora opcional
	\vspace{0.5cm}
	
	% --- Transcrição ---
	\section*{2. Transcrição / Notas da Entrevista}
	
	\begin{transcription}
		
		\pergunta{Quando um cliente chega com uma trotinete, quais são as primeiras 3 informações que lhe pede obrigatoriamente?}
		\resposta{Quando um cliente chega com uma trotinete, as três primeiras informações que peço são o nome completo, o contacto telefónico e uma descrição do problema que a trotinete tem. Se for um cliente novo, também peço o número de contribuinte e, se possível, o e-mail para podermos emitir a fatura corretamente e manter o registo atualizado. Se for um cliente habitual, tento confirmar se os dados estão corretos ou se houve alguma alteração.}
		
		\pergunta{É comum o cliente já estar registado? Gostaria que o sistema encontrasse o cliente apenas pelo número de telemóvel ou pela matrícula/ID da trotinete?}
		\resposta{É bastante comum o cliente já estar registado, principalmente os que vêm com frequência ou empresas que têm várias trotinetes. Por isso, seria muito útil se o sistema conseguisse encontrar o cliente apenas com o número de telemóvel ou até com a matrícula ou número de série da trotinete. Isso pouparia tempo e evitaria duplicações de registos, que às vezes acontecem quando o cliente não se lembra se já veio cá antes.}
		
		\pergunta{Como é que regista o 'estado visual' da trotinete à entrada (riscos, peças partidas) para evitar reclamações mais tarde? Escreve num campo de notas ou gostava de anexar fotos?}
		\resposta{Quanto ao estado visual da trotinete, costumo escrever à mão ou num campo de observações se vejo riscos, peças partidas ou sinais de queda. Tento ser o mais detalhada possível, mas nem sempre é fácil lembrar-me de tudo. Já tivemos situações em que o cliente reclamou de um risco que já lá estava, por isso seria ótimo se pudéssemos tirar uma ou duas fotos no momento da entrada e anexar ao registo. Assim ficava tudo documentado e evitavam-se mal-entendidos.}
		
		\pergunta{Como descreve o problema que o cliente relata? Tem uma lista de avarias comuns (ex: 'furo', 'não liga') ou escreve sempre um texto livre?}
		\resposta{Quando o cliente explica o problema, escrevo o que ele me diz com as minhas palavras, num campo de texto livre. Às vezes uso expressões como “não liga”, “faz barulho”, “não trava bem” ou “a bateria não dura”, mas não tenho uma lista fixa. Se houvesse uma lista com os problemas mais comuns, ajudava a uniformizar os registos e a facilitar a vida dos mecânicos, mas também acho importante poder escrever à parte se o cliente disser algo mais específico ou fora do normal.}
		
		\pergunta{Como é que sabe que data de entrega deve prometer ao cliente? O sistema deveria mostrar-lhe quantas trotinetes já estão na fila de espera?}
		\resposta{Quando um cliente chega com uma trotinete, costumo perguntar logo se tem alguma urgência, porque isso ajuda-me a perceber se posso prometer uma entrega mais rápida ou se tenho de explicar que há outras à frente. Normalmente, falo com o mecânico ou com o gerente para perceber quantas estão na fila e quanto tempo estimam para cada uma. Não temos uma forma exata de ver isso, é mais por experiência e conversa. Se o sistema mostrasse quantas trotinetes estão em espera e em que fase estão, ajudava-me muito a dar uma previsão mais realista ao cliente.}
		
		\pergunta{Depois de criar a ordem de serviço, imprime algum comprovativo para dar ao cliente? O que é que tem de constar obrigatoriamente nesse papel?}
		\resposta{Depois de criar a ordem de serviço, costumo escrever um papel à mão ou imprimir um formulário simples com os dados principais: nome do cliente, contacto, modelo da trotinete, descrição do problema e data de entrada. Também coloco um número de ordem para facilitar uma procura posterior. O cliente leva esse papel como comprovativo e nós ficamos com uma cópia. Era importante que esse documento tivesse pelo menos o nome do cliente, o número da ordem, a descrição do problema e a data prevista de entrega.}
		
		\pergunta{Se o cliente ligar a perguntar pelo estado da reparação, onde é que consulta essa informação atualmente? É fácil de encontrar?}
		\resposta{Quando um cliente liga a perguntar pelo estado da reparação, vou procurar no caderno ou nas folhas onde temos as ordens de serviço. Às vezes tenho de ir perguntar ao mecânico ou ao gerente, porque nem sempre está tudo escrito. Não é muito prático, principalmente se estiver sozinha na receção e com mais clientes à espera. Se o sistema tivesse essa informação acessível por nome ou número da ordem, seria muito mais rápido.}
		
		\pergunta{Como é que avisa o cliente que a trotinete está pronta ou que o orçamento foi alterado? Gostaria que o sistema enviasse um e-mail ou SMS automático para lhe poupar tempo de chamadas?}
		\resposta{Para avisar o cliente de que a trotinete está pronta ou que o orçamento mudou, costumo ligar por telefone. Às vezes mando mensagem, mas nem sempre tenho tempo para isso. Se o sistema pudesse enviar um SMS ou e-mail automático com essa informação, poupava-me muito trabalho e evitava esquecimentos, principalmente nos dias mais cheios.}
		
		\pergunta{Quando o mecânico deteta um problema extra, como é que essa informação chega até si para que possa falar com o cliente?}
		\resposta{Quando o mecânico encontra um problema extra, normalmente vem até à receção e explica-me o que encontrou. Depois sou eu que ligo ao cliente para explicar e perguntar se quer avançar com a reparação. Mas às vezes, se estou ocupada, o mecânico acaba por ligar ele próprio. Se houvesse uma forma de ele registar isso no sistema e eu receber logo um aviso, ajudava a manter tudo mais organizado e a garantir que o cliente é contactado rapidamente.}
		
		\pergunta{No momento da entrega, como é que confirma que o cliente pagou tudo o que devia? O sistema deve bloquear a entrega se houver valores pendentes?}
		\resposta{No momento da entrega, confirmo se o cliente pagou tudo o que devia consultando o talão da fatura ou o registo que temos no caderno ou no programa de faturação. Se não houver registo de pagamento, não entregamos a trotinete. Já aconteceu algumas vezes o cliente dizer que depois paga, mas isso dá problemas, por isso agora só entregamos com tudo pago. Seria ótimo que o sistema bloqueasse automaticamente a entrega se ainda houvesse valores em dívida, para evitar esquecimentos ou enganos.}
		
		\pergunta{Quais são as formas de pagamento que mais usa (numerário, multibanco, MBWay)? O sistema precisa de registar qual foi o método usado em cada venda?}
		\resposta{As formas de pagamento mais usadas são numerário, multibanco e MBWay. Cada cliente tem a sua preferência, mas ultimamente o MBWay tem sido muito comum, especialmente entre os mais jovens. Sim, era importante que o sistema registasse qual foi o método usado em cada venda, até para depois sabermos quanto entrou em dinheiro, quanto foi por cartão, e assim conseguimos fechar o dia com mais controlo.}
		
		\pergunta{É comum os clientes pedirem fatura com número de contribuinte à última da hora? É fácil para si alterar esses dados no momento de fechar o serviço?}
		\resposta{É muito comum os clientes pedirem fatura com número de contribuinte mesmo no fim, quando já está tudo pronto. Às vezes até já imprimimos a fatura e depois temos de voltar atrás para corrigir. Por isso, era importante que fosse fácil alterar esses dados no momento de fechar o serviço, sem ter de refazer tudo. Se o sistema permitisse editar só o campo do contribuinte antes de finalizar, ajudava imenso.}
		
		\pergunta{Há alguma tarefa que faça muitas vezes ao dia e que sinta que perde muito tempo (ex: preencher os mesmos dados várias vezes)? Como é que o sistema poderia tornar isso num clique?}
		\resposta{Uma tarefa que faço muitas vezes ao dia e que me faz perder tempo é repetir os mesmos dados em vários sítios. Por exemplo, quando crio uma ordem de serviço, tenho de escrever o nome do cliente, o contacto, o modelo da trotinete, e depois, na hora de faturar, tenho de voltar a escrever tudo outra vez. Se o sistema puxasse automaticamente esses dados, bastava confirmar e seguir em frente. Isso poupava muito tempo, principalmente quando temos vários clientes à espera.}
		
		\pergunta{Se se enganar a registar uma entrada, prefere conseguir apagar tudo ou prefere que o sistema a deixe apenas corrigir o erro?}
		\resposta{Se me enganar a registar uma entrada, prefiro que o sistema me deixe corrigir o erro em vez de apagar tudo. Às vezes é só um número trocado ou um nome mal escrito, e não faz sentido ter de apagar e começar de novo. Mas também acho que devia haver um histórico das alterações, para sabermos quem mudou o quê, caso seja preciso verificar mais tarde.}
		
		\vspace{0.8cm}
		
		\hrule
			
		\vspace{0.8cm}
		
		\pergunta{Quando tem uma fila de 3 clientes, qual é o tempo máximo que o sistema pode demorar a abrir uma ficha de cliente para não gerar impaciência?}
		\resposta{Se tenho três clientes à espera, o sistema não pode demorar mais do que 2 ou 3 segundos a abrir uma ficha de cliente. Se demorar mais do que isso, começo a sentir que estou a atrasar o atendimento e os clientes ficam impacientes. Tudo o que for lento ou exigir muitos cliques seguidos atrapalha bastante nesses momentos de maior pressão.}
		
		\pergunta{Se estiver a meio de um registo e o telefone tocar, o sistema permite-lhe 'minimizar' essa tarefa e abrir outra sem perder o que já escreveu?}
		\resposta{Se estiver a meio de um registo e o telefone tocar, o ideal é que o sistema me permita minimizar essa tarefa e abrir outra janela ou separador, sem perder o que já escrevi. Muitas vezes tenho de interromper o que estou a fazer para atender chamadas ou procurar outra informação, e se o sistema apagar os dados ou não me deixar voltar ao ponto onde estava, é muito frustrante.}
			
		\pergunta{Ao procurar por um cliente, basta-lhe o nome ou o sistema tem de ser capaz de encontrar resultados instantaneamente através do telemóvel ou NIF?}
		\resposta{Quando procuro um cliente, o nome nem sempre é suficiente — há muitos clientes com nomes parecidos ou que não se lembram do nome com que se registaram. Por isso, o sistema tem mesmo de conseguir encontrar resultados instantaneamente através do número de telemóvel ou do NIF. Esses dados são únicos e ajudam a evitar confusões.}
		
		\pergunta{O ecrã do computador fica visível para os clientes que estão ao balcão? Precisamos de uma opção para 'esconder' valores de faturação ou dados de outros clientes rapidamente?}
		\resposta{O ecrã do computador fica parcialmente visível para os clientes que estão ao balcão, especialmente se se inclinarem ou se estiverem do lado. Por isso, seria muito importante ter uma opção rápida para esconder valores de faturação ou dados de outros clientes — por exemplo, um botão de “modo privado” ou uma vista simplificada só com o essencial.}
		
		\pergunta{Como lidamos com o RGPD (Proteção de Dados)? O sistema deve pedir uma confirmação visual de que o cliente aceita receber SMS/e-mails de marketing?}
		\resposta{Em relação ao RGPD, temos de ter cuidado com os dados dos clientes. O sistema devia pedir uma confirmação visual — como uma caixinha de verificação — para sabermos se o cliente aceita ou não receber SMS ou e-mails de marketing. Assim ficava tudo registado e evitávamos problemas legais.}
		
		\pergunta{Se precisar de se ausentar do balcão por 2 minutos, o sistema deve bloquear-se automaticamente ou prefere bloquear manualmente para evitar que estranhos mexam no PC?}
		\resposta{Se precisar de me ausentar do balcão por 2 minutos, prefiro que o sistema se bloqueie automaticamente após um curto período de inatividade. Às vezes saio para ir buscar uma trotinete ou falar com um mecânico, e não dá tempo de bloquear manualmente. Mas também devia haver um botão de bloqueio rápido, para usar quando sei que vou sair por mais tempo.}
	
	\end{transcription}

%	\vspace{0.8cm}
%	
%	% --- Observações e Próximos Passos ---
%	\section*{3. Análise e Observações Gerais}
%	\begin{itemize}[leftmargin=*]
%		\item \textbf{Tom de voz:} O entrevistado demonstrou abertura, mas alguma frustração com os sistemas de IT atuais.
%		\item \textbf{Pontos Críticos:} Falta de interoperabilidade entre sistemas.
%	\end{itemize}
%	
%	\section*{4. Próximos Passos}
%	\begin{tabularx}{\textwidth}{@{}|X|l|l|@{}}
%		\hline
%		\rowcolor{gray!10} \textbf{Ação} & \textbf{Responsável} & \textbf{Prazo} \\ \hline
%		Validar fluxo de dados com equipa de IT & [Teu Nome] & 05/03/2026 \\ \hline
%		Agendar follow-up para demonstração de solução & [Teu Nome] & 12/03/2026 \\ \hline
%	\end{tabularx}
	
\end{document}